\chapter{Detailed Related-Work Boundary Comparison}
\label{app:related-work-boundaries}

Table~\ref{tab:background-related-work} in Chapter~\ref{chap:background} compares protected assets, attacker capabilities, and observables. The complementary table below retains the detailed policy, state, and evaluation-boundary comparison.

{\footnotesize
\setlength{\tabcolsep}{3pt}
\renewcommand{\arraystretch}{1.08}
\begin{longtable}{@{}>{\raggedright\arraybackslash}p{22mm}>{\raggedright\arraybackslash}p{39mm}>{\raggedright\arraybackslash}p{35mm}>{\raggedright\arraybackslash}p{40mm}@{}}
\caption{Policy scope, state, and evaluation boundaries in directly relevant work.}
\label{tab:background-related-work-boundaries}\\
\toprule
Work & Structure and policy granularity & State or authority transition & Control or evaluation boundary \\
\midrule
\endfirsthead
\multicolumn{4}{l}{\textit{Table~\ref{tab:background-related-work-boundaries} continued}}\\
\toprule
Work & Structure and policy granularity & State or authority transition & Control or evaluation boundary \\
\midrule
\endhead
\midrule
\multicolumn{4}{r}{\textit{Continued on next page}}\\
\endfoot
\bottomrule
\endlastfoot

\citet{jeong2025permissionrag} & Native IAM checks at resource/document level; field and relation policy are not central evaluated features & Authority transition not evaluated or reported & Post-retrieval permission filtering, access correctness, answer quality, and latency; an application raw/delivery boundary is not a central reported feature \\
\addlinespace[1.5pt]
\citet{klisura2026role} & Table- and column-level PostgreSQL policies; Text-to-SQL rather than RAG & Roles vary across examples; retained-state downgrade is outside the evaluated scope & Prompting, fine-tuning, and generator--verifier comparison; RAG-stage telemetry is outside scope \\
\addlinespace[1.5pt]
\citet{fazlija2025accessdenied} & Attribute-level policy over synthetic corporate records; independent relation-edge policy is not evaluated & Single exchange; no retained-state authority transition & Self-contained prompt with the relevant data and rules; output grading without operational retrieval stages \\
\addlinespace[1.5pt]
\citet{fazlija2026sensitivity} & RBAC-derived policy over sensitive and non-sensitive employee attributes; independent relation-edge policy is not evaluated & Per-session formal and empirical evaluation; no retained high-to-low transition & Formal RAG oracle and sensitivity-awareness privacy game; empirical LoRA and output comparison without matched application-guard attribution \\
\addlinespace[1.5pt]
\citet{zeng2024ragprivacy} & Retrieval records; role-conditioned field/relation authorisation is not evaluated & Authority transition not evaluated or reported & End-to-end extraction and privacy analysis; an application raw/delivery guard boundary is not reported \\
\addlinespace[1.5pt]
\citet{qi2025spill} & Text datastore; role-conditioned field/relation policy is outside the evaluated scope & Multiple queries, not a retained-authority transition & Extraction through instruction following; evaluation centres on generated output rather than an application delivery guard \\
\addlinespace[1.5pt]
\citet{anderson2025ragmia} & Passage membership; sensitivity-aware record fields and edges are outside the evaluated scope & -- & Membership attack evaluation at the query/response boundary; application-stage telemetry is not a central feature \\
\addlinespace[1.5pt]
\citet{choi2025safeguarding} & Retrieval-record membership rather than role-governed fields or edges & -- & Similarity-based query detection and detect-and-hide defence across retrieval and query/response boundaries \\
\addlinespace[1.5pt]
\citet{liu2026graphragprivacy} & Explicit graph structure; role-conditioned field/edge policy is not a central evaluated feature & -- & End-to-end extraction and candidate defences; thesis-style state/raw/delivery separation is not reported \\
\addlinespace[1.5pt]
\citet{greshake2023indirect} & External-content trust rather than structured field/relation authorisation & Application-dependent & Establishes indirect prompt injection across integrations; evaluation centres on attacks rather than matched RAG guard attribution \\
\addlinespace[1.5pt]
\citet{zhong2023retrievalpoisoning} & Text passages; field/relation confidentiality policy is outside the evaluated scope & -- & Retrieval manipulation under specified attacker capabilities and selected defence conditions \\
\addlinespace[1.5pt]
\citet{zou2025poisonedrag} & Text passages; structured authorisation is outside the evaluated scope & -- & Knowledge-corruption attack and selected defence evaluations at retrieval/generation boundaries \\
\addlinespace[1.5pt]
\citet{chaudhari2026phantom} & Text passages; role-conditioned field/relation policy is outside the evaluated scope & Trigger condition, not an authority transition & Retriever/generator optimisation, transfer tests, and end-to-end black-box evaluation; application policy and delivery attribution are outside scope \\
\addlinespace[1.5pt]
\citet{chen2024agentpoison} & Retrieved examples or knowledge; role-conditioned field/relation authorisation is outside scope & Long-term memory is an attack surface; legitimate role downgrade is not evaluated & External-memory/knowledge-base backdoor evaluation with attack success and benign performance \\
\addlinespace[1.5pt]
\citet{rao2026fragfuse} & Memory fragments; structured field/relation policy is outside the evaluated scope & Explicit multi-turn memory use; authority remains unprivileged & End-to-end bypass across access controls; a high-to-low role transition is not evaluated \\
\addlinespace[1.5pt]
\citet{liang2025saferag} & External text; role-conditioned field/relation confidentiality is outside scope & -- & Benchmark across fourteen components; raw-versus-delivered confidentiality and matched guard attribution are not central reported features \\
\addlinespace[1.5pt]
This thesis & Field sensitivity and independent relation policy over mixed-sensitivity entities & Legitimate higher-authority access followed by an explicit lower-authority transition & Stage-aware evaluation; package-level comparison is separated from component attribution \\

\end{longtable}
}
